% Fragmento para pegar en la sección de Resultados y análisis del Plan de Pruebas.
% Actualizar fecha, hora, equipo y número final de pruebas según la ejecución real.

\subsection{Resultados de pruebas Backend/API v0.1}

La primera versión de pruebas automatizadas se enfocó en los servicios de autenticación,
consulta del área de estudio y controles básicos de seguridad de la API. La ejecución se
realizó en ambiente local, sobre la rama \texttt{main}, utilizando una base SQLite temporal
para evitar modificaciones sobre la base de datos funcional del prototipo.

\begin{table}[H]
\centering
\begin{tabular}{|p{3cm}|p{7cm}|p{3cm}|}
\hline
\textbf{Identificador} & \textbf{Descripción} & \textbf{Estado} \\
\hline
API-01 & Registro exitoso de usuario. & Pendiente de ejecución \\
\hline
API-02 & Rechazo de usuario duplicado. & Pendiente de ejecución \\
\hline
API-03 & Rechazo de campos inválidos en registro. & Pendiente de ejecución \\
\hline
API-04 & Inicio de sesión exitoso y sesión autenticada. & Pendiente de ejecución \\
\hline
API-05 & Rechazo de contraseña incorrecta. & Pendiente de ejecución \\
\hline
API-06 & Consulta de sesión sin autenticación. & Pendiente de ejecución \\
\hline
API-07 & Cierre de sesión. & Pendiente de ejecución \\
\hline
API-08 & Consulta del área de estudio como GeoJSON. & Pendiente de ejecución \\
\hline
SEC-01 & Rechazo de exportación sin sesión. & Pendiente de ejecución \\
\hline
SEC-02 & Rechazo de consulta de estado sin sesión. & Pendiente de ejecución \\
\hline
SEC-03 & Rechazo de cancelación de exportación sin sesión. & Pendiente de ejecución \\
\hline
SEC-04 & Manejo controlado de archivo inexistente en exportaciones. & Pendiente de ejecución \\
\hline
SEC-05 & Validación de acceso indebido o ruta inválida en exportaciones. & Pendiente de ejecución \\
\hline
\end{tabular}
\caption{Resultados de pruebas Backend/API v0.1}
\end{table}

Como evidencia de ejecución se generó un archivo de log en
\texttt{tests\_evidence/backend\_api/}, el cual registra la salida completa de \texttt{pytest}.
